\section{Discussion} \label{sec:discussion}

This thesis investigated the users' mental models of \ac{LLM}-based systems in collaborative problem-solving scenarios, and the effect of explanations on those mental models. It derived design principles for \ac{XAI} in \ac{GenAI} systems to support accurate mental models of \ac{LLM}-based systems. The research questions were answered through a mixed-methods study, which combined qualitative process coding with quantitative survey data. The design principles were evaluated through expert interviews.

\subsection{Discussion of Results} \label{ssec:discussion-results}

% This paragraph is preseverd to be used to formulate a new introduction to the discussion once the mental model reframing is done.
% The attribution of mental states to an interactive system is spontaneous and largely unavoidable. Users apply the same folk-psychological framework they use to explain human behaviour to artificial agents \parencite{Miller2019}, a disposition that requires neither language nor embodiment \parencite{Heider1944} and that extends to autonomous and intelligent systems in general \parencite{deGraaf2017, Muller2004}. Natural-language interaction creates particularly favourable conditions for it \parencite{Weisz2025}. This is what makes \ac{ToM} the route to the user's mental model of the system: users build that model \textit{by} attributing mental states to the system. The attributions are what is observable; the model is inferred from them. The results of this study add two qualifications to this picture. Attribution is not only voiced but enacted, and it is selective in what it covers.

The evaluation of the study confirmed that the attribution of mental states to \ac{GenAI} systems is spontaneous and largely unavoidable as discussed in Section~\ref{ssec:tom-attributions}.
% TODO: Insert reminder about ToM as lens for observing mental models here once the reframing is done.
% Note: The middle part is missing making writing this introductory paragraph quite hard. This will likely need a rewrtite.
Through the development of the design theory, this thesis surfaces several key findings which are discussed in the following. The discussion closes with clear answers to each of the research questions.

Participants attributed beliefs, desires, and intentions to the chatbot, but not uniformly across the three components. This is consistent with the findings that people use the same conceptual framework to explain the behaviour of artificial agents as they do for humans \parencite{deGraaf2017,Miller2019}. \textit{Beliefs} were attributed most explicitly, concentrated in the \textit{Mental Model \& Updates} dimension. Attributions typically surfaced when participants perceived a discrepancy between the system's assumptions and the information given in the task. \textit{Desires} were often attributed indirectly, through the participants' repetition and uptake of the system's responses. \textit{Intentions} were rarely attributed. Of the four signatures of meta-representation \parencite{Rakoczy2022}, participants utilised three regularly. \textit{Difference} and \textit{incompatibility} were used to establish a conflict between the participant's and the system's representation of the task. \textit{Misrepresentation} was reached when participants decided to correct the system, and specifically when they traced the system's misrepresentation back to their own input (\enquote{attributing LLM's confusion to own action}). \textit{Aspectuality} was not reached.

As a consequence of the indirect nature of attribution, attribution happens continuously across the interaction cycle rather than at discrete points. The three steps of the interaction cycle are represented by the three \ac{LLM}-directed aggregate dimensions: Users decide whether and how to engage the system (\textit{LLM Reliance Decisions}), exchange information with it (\textit{Information Exchange}), and appraise what they receive (\textit{Mental Model \& Updates}). The decision step mostly carries implicit attributions, as the decision to delegate is the enactment of a belief about system capability. Attribution during information exchange is implicit as well surfacing through what the user communicates to the system. The appraisal step carries direct attributions as well as implicit ones, through evaluation and uptake of the system's responses.

The appraisal step shows a clear default in favour of the system. While the themes and dimensions constructed do not allow cleanly separating the two types of conflict described by \textcite{Jussupow2021} without additional coding, the appraisal imbalance reported in Section~\ref{ssec:initial-design-principles} implies a default attribution of competence to the system that is not necessarily rooted in facts. Participants' scrutiny of their own work converted into a correction at nearly five times the rate it did for the system's output. An alternative interpretation is that the system simply erred less often than the participants. While this cannot be refuted with the data collected, the design implications are the same, especially reflected in \textbf{MR2}. A default attribution of competence to the system and resulting reliance on it would have been appropriate by accident rather than driven by validation. With a less competent system, the same default would have led to over-reliance (cf. \texttt{08.048.6}).

Where the collaboration issues arose, the problem most often traced back to the \textit{knowledge distribution} dimension of the users mental model \parencite{Gero2020}: participants held a misrepresentation of what the system knows. Two \ac{DP} seek to address this challenge directly: \textbf{DP3} and \textbf{DP5} improve the communication of the system's knowledge distribution to the user. Beyond that, \textbf{DP2} and \textbf{DP-T5} seek to support the users' model of \textit{global behaviour}. \textit{Local Behaviour} understanding is support by \textbf{DP-T1} completing coverage of categories of mental model content by \parencite{Gero2020}.

% TODO: Isolate this as a contribution
% Beyond these, \textbf{DP7} and \textbf{DP-T3} invert the direction of interest, from the user's model of the system to the system's model of the user. This is precisely what \ac{MToM} advocates \parencite{Wang2021, Wang2024a}, and it constitutes an extension of \textcite{Gero2020}'s categories, which are formulated unidirectionally: the same three categories can be applied in the opposite direction.

The transcripts also show that participants use anthropomorphic vocabulary when referring to the chatbot. This behaviour is marked predominantly by volitional and knowledge-related verbs, e.g., \enquote{er will die ganze Zeit, dass ich \ldots} [he wants me to \ldots all the time] (ID: \texttt{03.020.2}) and \enquote{das kann er ja gar nicht wissen} [he cannot possibly know that] (ID: \texttt{11.022.4}). Rarely participants also exchanged social acts with the chatbot that would typically be associated with inter-human interaction, e.g., \enquote{Vielen Dank, Chatbot} [thank you, chatbot] (ID: \texttt{02.025.10}). While the anthropomorphic behaviour occurs across all participants, the density of these statements varies considerably. Notably, for some participants it alternates with agent framing, e.g., \enquote{es scheint die KI dann zu hängen} [the AI then seems to hang] (ID: \texttt{08.034.5}). The coding methodology applied does not offer data to further investigate these claims in a structured manner and German grammatical gender makes quantitative analysis based on pronouns infeasible. The evidence therefore remains anecdotal.

% TODO: This can likely just be dropped. It repeats the results chapter and adds no discussion of the greater research context
% The quantitative layer serves as triangulation only. \textbf{Frustration} ($p=.005$) and \textbf{Perceived Usefulness} ($p=.095$) differ significantly between the groups at $\alpha=0.1$, and the remaining workload constructs are directionally consistent with the hypothesised effect (cf. Section~\ref{sssec:gap-analysis}).

Referring back to the research questions posed in Section~\ref{ssec:research-objectives}, \textbf{RQ1} asks \enquote{which mental states \ac{LLM}-users attribute} and \enquote{their mental models of the system}. This thesis shows that beliefs and desires are attributed, and that attribution is to a significant degree implicit, enacted through behaviour rather than voiced. Explicit attributions surfaced mainly where participants encountered a discrepancy between the system's assumptions and the task. What can be said about the mental models themselves is a reconstruction from that attribution behaviour.

\textbf{RQ2} focused on the effects of \enquote{explanations for \ac{LLM} outputs} on mental model attribution. The data produced during the experiment does not allow for dimension- or theme-level between group comparisons with Section~\ref{sssec:gap-analysis} discussed potential reasons. The sole finding showed, that explanations decreased the share of heuristic uptake of \ac{LLM} outputs from 36.0\% to 17.4\%.

The last research question focusing on design knowledge derived from \textbf{RQ1} and \textbf{RQ2} is answered by the design theory developed throughout this thesis. Figure~\ref{fig:mr-dr-dp-mapping-revised} shows the expert evaluated final version of the design theory and Section~\ref{sssec:revised-dp} elaborates the \ac{DP} in a structured format. To address the gaps created by the limited answer to \textbf{RQ2} Section~\ref{sssec:gap-analysis} also presents theory-derived \ac{DP}, that have not been expert evaluated.

% TODO: Consider if this sentence might be something better placed in the conclusion or even an abstract, if shortened further?
% Attributions to the system were largely enacted rather than voiced: participants acted on an understanding of the system they did not articulate, and explicit attributions surfaced mainly where the system's behaviour conflicted with their expectations. In this study, such discrepancies were resolved in the system's favour almost five times more often than against it. Consequently, the design goal is not only to help users form an accurate mental model, but also to expose the existing one to challenge in situations where it would otherwise go unexamined.

\subsection{Contributions} \label{ssec:contributions}

The first contribution of this thesis is the proposition that ToM is a route through which users form and revise their mental models of an AI system (cf. Section~\ref{sssec:theory-of-mind-ai}). Prior work has shown that users can adopt mentalistic stances towards artificial agents \parencite{PerezOsorio2020}; this thesis proposes how such stances relate to the mental models users hold, and adopts that relation as the interpretive frame for the present analysis. In doing so it articulates one relation that \textcite{Bodamer2026} identify as characteristically left implicit in \ac{HAII} research. Secondly, the literature review in Section~\ref{sec:problem-awareness} structures recent relevant literature and organises it into three research directions and explaining their interrelations (summarised in Figure~\ref{fig:hai-genai-roadmap}).

The design theory developed in this thesis is itself the primary contribution to the field of \ac{HAII}. It builds on the design knowledge of \textcite{Amershi2019} and \textcite{Weisz2025}, carrying the user needs established by the former into the \ac{GenAI} setting and enriching the latter with a focus on the user's mental model. Beyond the \acp{MR}, \acp{DR}, and \acp{DP} stated in the previous chapters, the discussion surfaced an additional design goal. \textbf{MR1}~--~\textbf{MR3} aim to support the formation of an accurate mental model, presupposing that the user is able and willing to do so. Section~\ref{ssec:explanations-mental-model-formation} already notes this constraint: \textcite{Jussupow2021} suggest that learning from a conflict is achieved only where the user reaches \textit{active consideration}, which cannot be assumed. The emergent goal is therefore to expose the model a user already holds to challenge at the points where it influences their decision. \textbf{DR4} already serves this goal as a means, and the \acp{DP} addressing the goal are exactly those serving \textbf{DR4} (\textbf{DP4}, \textbf{DP6}, \textbf{DP-T3}, and \textbf{DP-T4}). What the \ac{MR} layer does not do is name it as an end: \textbf{MR3} comes closest, but is scoped to the mode in which a user processes a system's output rather than to the examination of the model through which that output is interpreted.

Another contribution is the identification of a failure mode that inverts a known one. \textit{AI-based confirmation} is a failure mode, in which the user anchors their belief on the system's advice before an independent frame has formed \parencite{Jussupow2021}. The data of this study show the inverse: the system adopts the user's erroneous input. The setting of \textcite{Jussupow2021}, using a discriminative \ac{CAID} system issuing binary advice, does not allow for this new failure mode, as the system cannot adopt a user's premise. While a single instance is sufficient to prove existence of this failure mode, its generalisation and its importance across contexts remain open. \textbf{DP4} is the counter-mechanism built to address this failure mode and two experts independently noted during the evaluation that surfacing conflicts works against the sycophantic tendencies of \acp{LLM}.

The final theoretical contribution are insights into the role of faithfulness in \ac{GenAI} design. The \ac{CoT} explanations used in this study are not guaranteed to reflect the process that produced an output \parencite{Lyu2023}. Because guaranteeing faithfulness remains unsolved \parencite{Tutek2025}, every \ac{DP} whose mechanism depends on the accuracy of an explanation or of a system self-report is exposed to this risk. The practical impact to for proposed \ac{DP} depends on the \ac{MR} and ultimately on the design goal it serves. The initial design goal of supporting the formation of an accurate mental model benefits is dependent on faithfulness and requires mitigation strategies. Conversely, the emergent design goal of exposing and challenging an existing mental model is less impacted: even an unfaithful explanation can surface an existing model and induce reflection. This means that \textbf{DP3}, \textbf{DP4}, \textbf{DP5}, and \textbf{DP-T1} are exposed to this risk, with \textbf{DP-T2} acting as explicit mitigation.

Three practical implications follow for practitioners. First, the availability of an explanation does not guarantee its consumption. The \ac{CoT} was rendered collapsed and click-to-expand (Section~\ref{sssec:model-selection}), and only two pattern codes touch it at all. However, this evidence is partial, as task difficulty and silent reading compete as explanations for the same observation (Section~\ref{sssec:gap-analysis}). Second, out-of-scope requests should be rejected clearly. The system was instructed not to solve the problems for users, even when directly prompted and users where briefed explicitly. A refusal should be paired with a disclosure of what the system can do (\textbf{DP2} and \textbf{DP-T5}), or it reads as a failure rather than as a boundary. Third, friction is a budget. \textbf{DR4} and \textbf{DR5} stand in deliberate tension (Section~\ref{ssec:meta-requirements}), and the experts raised habituation as a concern for \textbf{DP6}, which applies equally to \textbf{DP-T4}. Friction should be spent only where \textbf{DR4}'s consequentiality test is met, that is where an action is irreversible, costly to correct, or potentially harmful \parencite[cf.][]{Buinca2021}.

\subsection{Limitations \& Future Work} \label{ssec:limitations-future-work}

Two key limitations contributed to the gap in answering \textbf{RQ2}. The first is executional, arising from the design of the study. The task difficulty turned out to be below the required threshold to induce a need for explanation, as evidenced by 32 of the 36 tasks being solved correctly. Between group differences were therefore underpowered to the degree where differences could not be reliably traced back to the manipulation. The second is structural, arising from the retrospective nature of the analysis. The study was originally designed to investigate cognition and metacognition of users, not their mental models or \ac{ToM} attributions. Therefore, no \ac{ToM} instruments were administered, like the scales used by \textcite{Colombatto2025} or any measure of behavioural reliance. Beyond that, the study was conducted in a single session per participant, so the \enquote{over time} category of \textcite{Amershi2019} is unobservable.

The analytic lens of \ac{ToM} attribution was also limited in its reach. While observing attribution behaviour from qualitative transcripts is a valid approach \parencite{deGraaf2019}, the process-focused coding scheme only surfaces attributions that impact participants' actions. The codebook used in this study can therefore not distinguish between an attribution that did not occur and one that occurred but left no trace. As a result, the uneven coverage of the three components of \ac{ToM} can be interpreted in two ways: either participants attributed beliefs, desires, and intentions unequally, or the coding scheme surfaced the attributions unevenly.

% The reach of the analytical lens is bounded in a second way. The BDI frame was applied to finished pattern codes rather than built into the coding scheme, which applied process and in vivo coding first \parencite{Saldana2015}. Coverage is therefore uneven: belief is reached well, desire only where it coincided with a participant's process, and intention barely at all. This is a question of analytic coverage, not of measurement validity, as no validated measure of attribution is claimed. The observability conditions given in Section~\ref{ssec:discussion-results} explain which attributions the lens can reach and which it cannot, so the boundary of the picture is known rather than unknown. Beyond this, the mental models themselves were never observed. What is reported is \textcite{Norman1983}'s fourth element: the researcher's reconstruction of the user's model from the attribution behaviour available.

% Two limitations arise from the execution of the study. The task difficulty fell below the threshold at which an explanation becomes necessary, as evidenced by 32 of 36 tasks being solved correctly, which made engaging with the explanation optional rather than necessary. This affects the statistical power of the comparison rather than the design concept, and was identifiable only after execution. Further, the manipulation is not isolable: instructing the model to produce a chain of thought changes both the content of its responses \parencite{Wei2022, Kojima2022} and whether a trace is displayed at all, so the between-group evidence speaks only weakly to the effect of \textit{displaying} an explanation.

The analysis was conducted by a single researcher, so no intercoder reliability, negotiated agreement, or member checking was possible, contrary to the team-based assumption underlying \textcite{Gioia2013}. This exposure extends across the coding, the application of the BDI frame, the construction of the data structure, the abductive derivation of the design principles, and the completion of the evaluation forms. The hermeneutic literature review carries the same exposure, in addition to its inherent lack of reproducibility.

From these limitations, several directions for future research emerge. First, the complete design theory developed in this thesis (including \textbf{DP-T1}~--~\textbf{DP-T5}) should be instantiated and evaluated. This evaluation should measure important \ac{ToM}-appropriate constructs, like the dimensions of \textcite{Colombatto2025}, behavioural reliance, and the enactment of attributions. Additional amendments to the research design could include longitudinal design, to observe the evolution of attributions over time. Lastly, the coding scheme should be amended to facilitate the BDI frame.

An entirely different direction for future work is to investigate the failure mode of the \ac{AI} accidentally adopting a user's erroneous premise. The present study shows that this failure mode exists, but its generalisation and importance across contexts remain open. Furthermore, the systematic investigation of additional failure modes in \ac{GenAI} systems is a promising avenue for future research and their impact on users' mental models and \ac{ToM} attributions. Additionally, the connection between the user's mental model and task performance could be investigated, as this study cannot reliably speak to that connection.
